% GENERATED FILE. Edit canonical lessons, then run npm run generate:books.
% canonical-source-hash: fnv1a64:c676e7e39a499c1c
% canonical-lessons: GU-C537-pasandgi, GU-C537-badle, GU-C537-manapsand, GU-C537-vinvu, GU-C537-rucvu

\chapter{Choosing}
\label{ch:gu-choosing}
\hlchaptermodality{537}

\begin{chapteropening}
\textbf{By the end of this chapter:} \emph{I can say what I would choose and what I would take instead.}

Complete the last lesson of chapter 537: I can say what I would choose and what I would take instead.
\end{chapteropening}

\section[\texorpdfstring{\gu{પસંદગી} (pasandgī)}{પસંદગી (pasandgī)}]{\gu{પસંદગી} (pasandgī) — a choice}

\label{lesson:GU-C537-pasandgi}



\begin{quote}
Before the new one: say the Gujarati for a birth chart, then the Gujarati for the future.
\end{quote}

\subsection*{You'll want to know: \gu{પસંદગી}}
\textbf{\gu{પસંદગી}} — \emph{pasandgī} — \textquotedblleft{}a choice\textquotedblright{}.

\begin{quote}
\gu{જો} \gu{પસંદગી} \gu{મારી} \gu{હોત} \gu{તો} \gu{હું} \gu{પર્વત} \gu{પર} \gu{જાત}.
\end{quote}

\begin{quote}
\emph{jo pasandgī mārī hot to hũ parvat par jāt.}
\end{quote}

\begin{quote}
If the choice were mine, I would go to the mountains.
\end{quote}

A feminine noun.

\begin{cousinweb}
\textbf{\gu{પસંદગી}} (\emph{pasandgī}) is usually traced to Persian \emph{pasand}, \textquotedblleft{}approved, liked\textquotedblright{}, with the ending \emph{-gī}; \textbf{\gu{પસંદ} \gu{કરવું}} is to choose.
\end{cousinweb}

\begin{grammarlens}[title={What you've built}]
One more word for saying what happens if.
\end{grammarlens}

\subsection*{Your turn}
\begin{itemize}
\raggedright
  \item \emph{Say it:} \emph{pasandgī}
  \item \emph{Say it:} \emph{pasandgī}, once more
  \item \emph{Say it:} \emph{jo pasandgī mārī hot to hũ parvat par jāt.} — If the choice were mine, I would go to the mountains.
  \item \emph{Recall:} say the Gujarati for a birth chart, then the Gujarati for the future, then say \emph{pasandgī} again
\end{itemize}

\begin{tcolorbox}[breakable,colback=teal!4,colframe=teal!35!black,title={Before you move on}]
What does \gu{પસંદગી} mean? (\textquotedblleft{}A choice\textquotedblright{}.) Say it once more.
\end{tcolorbox}
\section[\texorpdfstring{\gu{બદલે} (badle)}{બદલે (badle)}]{\gu{બદલે} (badle) — instead of}

\label{lesson:GU-C537-badle}



\begin{quote}
Before the new one: say the Gujarati for the future, then the Gujarati for a choice.
\end{quote}

\subsection*{You'll want to know: \gu{બદલે}}
\textbf{\gu{બદલે}} — \emph{badle} — \textquotedblleft{}instead of\textquotedblright{}.

\begin{quote}
\gu{જો} \gu{ચા} \gu{નહીં} \gu{હોય} \gu{તો} \gu{હું} \gu{ચાને} \gu{બદલે} \gu{દૂધ} \gu{પીશ}.
\end{quote}

\begin{quote}
\emph{jo cā nahī̃ hoy to hũ cāne badle dūdh pīś.}
\end{quote}

\begin{quote}
If there's no tea, I'll drink milk instead.
\end{quote}

\begin{cousinweb}
\textbf{\gu{બદલે}} (\emph{badle}) is usually traced to Arabic \emph{badal}, \textquotedblleft{}a substitute\textquotedblright{}, through Persian, in the form that follows \textbf{\gu{ને}}: \textbf{\gu{ચાને} \gu{બદલે}}, instead of tea. The verb \textbf{\gu{બદલવું}}, to change, belongs with it.
\end{cousinweb}

\begin{grammarlens}[title={What you've built}]
One more word for saying what happens if.
\end{grammarlens}

\subsection*{Your turn}
\begin{itemize}
\raggedright
  \item \emph{Say it:} \emph{badle}
  \item \emph{Say it:} \emph{badle}, once more
  \item \emph{Say it:} \emph{jo cā nahī̃ hoy to hũ cāne badle dūdh pīś.} — If there's no tea, I'll drink milk instead.
  \item \emph{Recall:} say the Gujarati for the future, then the Gujarati for a choice, then say \emph{badle} again
\end{itemize}

\begin{tcolorbox}[breakable,colback=teal!4,colframe=teal!35!black,title={Before you move on}]
What does \gu{બદલે} mean? (\textquotedblleft{}Instead of\textquotedblright{}.) Say it once more.
\end{tcolorbox}
\section[\texorpdfstring{\gu{મનપસંદ} (manapsand)}{મનપસંદ (manapsand)}]{\gu{મનપસંદ} (manapsand) — favourite}

\label{lesson:GU-C537-manapsand}



\begin{quote}
Before the new one: say the Gujarati for a choice, then the Gujarati for instead of.
\end{quote}

\subsection*{You'll want to know: \gu{મનપસંદ}}
\textbf{\gu{મનપસંદ}} — \emph{manapsand} — \textquotedblleft{}favourite\textquotedblright{}.

\begin{quote}
\gu{જો} \gu{દુકાનમાં} \gu{મારો} \gu{મનપસંદ} \gu{રંગ} \gu{હશે} \gu{તો} \gu{હું} \gu{તે} \gu{સાડી} \gu{ખરીદીશ}.
\end{quote}

\begin{quote}
\emph{jo dukānmā̃ māro manapsand rang haśe to hũ te sāḍī kharīdīś.}
\end{quote}

\begin{quote}
If the shop has my favourite colour, I will buy that sari.
\end{quote}

\begin{cousinweb}
\textbf{\gu{મનપસંદ}} (\emph{manapsand}) is built from \emph{man}, \textquotedblleft{}the mind\textquotedblright{}, from Sanskrit \emph{manas} (Monier-Williams), and \emph{pasand}, \textquotedblleft{}liked\textquotedblright{}, usually traced to Persian \emph{pasand}: liked by one's own mind.
\end{cousinweb}

\begin{grammarlens}[title={What you've built}]
One more word for saying what happens if.
\end{grammarlens}

\subsection*{Your turn}
\begin{itemize}
\raggedright
  \item \emph{Say it:} \emph{manapsand}
  \item \emph{Say it:} \emph{manapsand}, once more
  \item \emph{Say it:} \emph{jo dukānmā̃ māro manapsand rang haśe to hũ te sāḍī kharīdīś.} — If the shop has my favourite colour, I will buy that sari.
  \item \emph{Recall:} say the Gujarati for a choice, then the Gujarati for instead of, then say \emph{manapsand} again
\end{itemize}

\begin{tcolorbox}[breakable,colback=teal!4,colframe=teal!35!black,title={Before you move on}]
What does \gu{મનપસંદ} mean? (\textquotedblleft{}Favourite\textquotedblright{}.) Say it once more.
\end{tcolorbox}
\section[\texorpdfstring{\gu{વીણવું} (vīṇvũ)}{વીણવું (vīṇvũ)}]{\gu{વીણવું} (vīṇvũ) — to pick out (one by one)}

\label{lesson:GU-C537-vinvu}



\begin{quote}
Before the new one: say the Gujarati for instead of, then the Gujarati for favourite.
\end{quote}

\subsection*{You'll want to know: \gu{વીણવું}}
\textbf{\gu{વીણવું}} — \emph{vīṇvũ} — \textquotedblleft{}to pick out (one by one)\textquotedblright{}.

\begin{quote}
\gu{જો} \gu{તમે} \gu{સારી} \gu{કેરીઓ} \gu{વીણશો} \gu{તો} \gu{હું} \gu{તે} \gu{ખરીદીશ}.
\end{quote}

\begin{quote}
\emph{jo tame sārī kerīo vīṇśo to hũ te kharīdīś.}
\end{quote}

\begin{quote}
If you pick out the good mangoes, I'll buy them.
\end{quote}

Past, with \gu{તેણે} (he or she …): \gu{તેણે} \gu{કેરીઓ} \gu{વીણી}, \emph{teṇe kerīo vīṇī}.

\begin{cousinweb}
\textbf{\gu{વીણવું}} (\emph{vīṇvũ}) is inherited from Sanskrit \emph{vícinōti}, \textquotedblleft{}picks out, sorts\textquotedblright{}, and CDIAL lists the Gujarati verb, \textquotedblleft{}to pick, gather\textquotedblright{} (CDIAL 11686).
\end{cousinweb}

\begin{grammarlens}[title={What you've built}]
One more word for saying what happens if.
\end{grammarlens}

\subsection*{Your turn}
\begin{itemize}
\raggedright
  \item \emph{Say it:} \emph{vīṇvũ}
  \item \emph{Say it:} \emph{vīṇvũ}, once more
  \item \emph{Say it:} \emph{jo tame sārī kerīo vīṇśo to hũ te kharīdīś.} — If you pick out the good mangoes, I'll buy them.
  \item \emph{Recall:} say the Gujarati for instead of, then the Gujarati for favourite, then say \emph{vīṇvũ} again
\end{itemize}

\begin{tcolorbox}[breakable,colback=teal!4,colframe=teal!35!black,title={Before you move on}]
What does \gu{વીણવું} mean? (\textquotedblleft{}To pick out (one by one)\textquotedblright{}.) Say it once more.
\end{tcolorbox}
\section[\texorpdfstring{\gu{રુચવું} (rucvũ)}{રુચવું (rucvũ)}]{\gu{રુચવું} (rucvũ) — to appeal (to someone)}

\label{lesson:GU-C537-rucvu}



\begin{quote}
Before the new one: say the Gujarati for favourite, then the Gujarati for to pick out.
\end{quote}

\subsection*{You'll want to know: \gu{રુચવું}}
\textbf{\gu{રુચવું}} — \emph{rucvũ} — \textquotedblleft{}to appeal (to someone)\textquotedblright{}.

\begin{quote}
\gu{જો} \gu{આ} \gu{યોજના} \gu{તમને} \gu{રુચશે} \gu{તો} \gu{આપણે} \gu{કાલે} \gu{શરૂ} \gu{કરીશું}.
\end{quote}

\begin{quote}
\emph{jo ā yojnā tamne rucśe to āpaṇe kāle śarū karīśũ.}
\end{quote}

\begin{quote}
If this plan appeals to you, we will start tomorrow.
\end{quote}

Past (it …): \gu{તે} \gu{વાત} \gu{તેને} \gu{રુચી}, \emph{te vāt tene rucī}.

\begin{cousinweb}
\textbf{\gu{રુચવું}} (\emph{rucvũ}) is inherited from Sanskrit \emph{rucyatē}, \textquotedblleft{}is pleasing\textquotedblright{}, and CDIAL lists the Gujarati verb, \textquotedblleft{}to be liked\textquotedblright{} (CDIAL 10765). The person it appeals to takes \textbf{\gu{ને}}.
\end{cousinweb}

\begin{grammarlens}[title={What you've built}]
One more word for saying what happens if.
\end{grammarlens}

\subsection*{Your turn}
\begin{itemize}
\raggedright
  \item \emph{Say it:} \emph{rucvũ}
  \item \emph{Say it:} \emph{rucvũ}, once more
  \item \emph{Say it:} \emph{jo ā yojnā tamne rucśe to āpaṇe kāle śarū karīśũ.} — If this plan appeals to you, we will start tomorrow.
  \item \emph{Recall:} say the Gujarati for favourite, then the Gujarati for to pick out, then say \emph{rucvũ} again
\end{itemize}

\begin{tcolorbox}[breakable,colback=teal!4,colframe=teal!35!black,title={Before you move on}]
What does \gu{રુચવું} mean? (\textquotedblleft{}To appeal (to someone)\textquotedblright{}.) Say it once more.
\end{tcolorbox}
